% ---
% id: meat-core/rem:scalar-declaration
% kind: rem
% title: "Scalar declaration of a homogeneous right-hand side"
% summary: >-
%   A trusted-mode entry declares the scalar c_k, not an operator Θ_k said to
%   equal c_k·τ_A, so its right-hand side and the transported bound are one
%   expression. An operator only near c_k·τ_A would need a numerical test, and
%   no relaxation of the program covers admitting it.
% prominence: standard
% status: accepted
% uses: [def:homog, lem:range]
% mentions: [def:feasible, meat-core/frank-wolfe/def:relaxed-feasible-set]
% assumes: []
% refs: []
% provenance: newly-derived
% lean:
% ---

\begin{remark}\label{rem:scalar-declaration}
Licensing is a type check in both modes and compares no numbers
(Definition~\ref{def:homog}). This is deliberate, and it is the reason the
trusted-mode entry declares ``a declared $c_k$'' and not ``a declared
$\Theta_k$ satisfying $\Theta_k=c_k\tau_A$''. Under the latter form, a
declaration carrying an operator $\Theta_k$ merely \emph{near} $c_k\tau_A$
would have to be admitted or refused by a numerical test, and admitting one
would be unsafe in a way no relaxation of this module covers: the transport
Lemma~\ref{lem:range}\,(i) and \eqref{eq:transport} deliver exactly
$c_k\tau_A$ on the right-hand side, while \eqref{eq:feas-modes} would impose
the declared $\Theta_k$, and the two feasible sets are then neither nested nor
related by the relaxation of Definition~\ref{def:relaxed-feasible-set}, whose
window is a spectral floor, or an entrywise ball, on the \emph{relation}, and
not an entrywise-relative window on the right-hand side. With the scalar
declaration, $\Theta_k$ and the transported bound are the same expression in
the same $c_k$ and the same $\tau_A$, and the two agree by construction on the
declared numbers.
\end{remark}
